% Part: history
% Chapter: biographies 
% Section: emmy-noether

\documentclass[../../../include/open-logic-section]{subfiles}

\begin{document}

\olfileid[tr]{his}{bio}{noe}

\olsection{Emmy Noether}

\olphoto{noether-emmy}{Emmy Noether}

Emmy Noether (\textsc{ner}-ter), 23 Mart 1882'de Almanya'nın Erlangen
kentinde, üst orta sınıfa mensup akademik bir ailede doğdu. ``Modern cebirin
annesi'' diye anılan Noether, kadınların eğitimine getirilen önemli engellere
karşın hem matematiğe hem fiziğe çığır açıcı katkılarda bulundu. Dönemin
Almanyası'nda genç kızların sanat alanında eğitim görmesi bekleniyor ve
üniversiteye hazırlık okullarına gitmelerine izin verilmiyordu. Bununla
birlikte Noether, G\"{o}ttingen ve Erlangen üniversitelerindeki dersleri
dinleyici olarak izledikten sonra (babası Erlangen'de matematik profesörüydü),
kuralların kadın öğrencilerin kabulüne izin verecek biçimde değiştirilmesiyle
1904'te Erlangen'e öğrenci olarak kaydolabildi. Matematik doktorasını 1907'de
aldı.

Noether, yeterliliğine rağmen kariyeri boyunca büyük dirençle karşılaştı.
1908--1915 arasında Erlangen'de ücretsiz ders verdi. Bu sırada dönemin
dünyaca önde gelen matematikçilerinden David Hilbert'in dikkatini çekti ve
Hilbert onu G\"{o}ttingen'e davet etti. Ne var ki kadınların profesörlük
alması yasaktı; bu nedenle yine ücret almadan, yalnızca Hilbert'in adı altında
ders verebildi. Bu dönemde, bugün Noether teoremi diye bilinen ve kuramsal
fizikte hâlâ kullanılan teoremi ispatladı. Noether'e nihayet 1919'da ders
verme hakkı tanındı. Üniversitedeki meslektaşlarının süren direncine Hilbert'in
şöyle karşılık verdiği aktarılır: ``Beyler, fakülte senatosu hamam değildir.''

1920'lerin sonlarına doğru soyut cebir çalışmalarına yoğunlaştı ve katkıları
bu alanda devrim yarattı. İspatlarında sık sık, belirli kümelerin sonsuz ve
kesin artan bir zincirinin bulunmadığını söyleyen sözde artan zincir koşulunu
kullandı. Örneğin günümüzde Noether halkaları denen bazı cebirsel yapılarda
$I_1 \subsetneq I_2 \subsetneq \dots$ biçiminde sonsuz ideal dizileri yoktur.
Bu koşul herhangi bir kısmi sıralamaya genelleştirilebilir (cebirde, altküme
bağıntısıyla sıralanan idealler özel durumunu ilgilendirir); ayrıca bir kısmi
sıralamadaki her kesin \emph{azalan} dizinin sonunda bittiği ikili azalan
zincir koşulunu da ele alabiliriz. Bir kısmi sıralama azalan zincir koşulunu
sağlıyorsa bu sıralama boyunca, $<$ ile sıralanmış~$\Nat$ boyunca tümevarım
kullanmamıza benzer biçimde tümevarım kullanmak mümkündür. Böyle
sıralamalara \emph{iyi temelli} ya da \emph{Noether tipi}, karşılık gelen
ispat ilkesine de \emph{Noether tümevarımı} denir.

Noether Yahudiydi ve Naziler 1933'te iktidara gelince görevinden çıkarıldı.
Neyse ki Amerika Birleşik Devletleri'ne göç ederek Pennsylvania'daki Bryn
Mawr'da geçici bir görev alabildi. Burada bulunduğu sırada Princeton'da da
ders verdi; ancak üniversiteyi kadınlara karşı misafirperver bulmadı
\citep[81]{Dick1981}. Noether 1935'te bir rahim tümörünün alınması için
ameliyat oldu. Ameliyatın yol açtığı bir enfeksiyon nedeniyle öldü ve Bryn
Mawr'a gömüldü.

\begin{reading} 
Noether'in bir yaşam öyküsü için \citet{Dick1981} kaynağına bakın. Perimeter
Kuramsal Fizik Enstitüsü'nün Noether'in yaşamı ve etkisi üzerine dersleri
çevrimiçi olarak mevcuttur \citep{Perimeter2015}. Okumaktan yorulduysanız
\emph{Stuff You Missed in History Class}, Noether'in yaşamı ve etkisi üzerine
bir podcast sunar \citep{Frey2015}. Noether'in toplu eserleri özgün
Almancalarıyla mevcuttur \citep{Noether1983}.
\end{reading}

\end{document}
